\section{Brain Topography: A Posterior-Visual Scope Statement}
\label{sec:topography}

The cohort signal is posterior-dominant: region-mean correlation is
strongest at occipital electrodes ($|r|=0.21$) and weakest at frontal
($|r|=0.16$), with the largest single cell at PO3/$\gamma$ ($r=+0.48$;
Figure~\ref{fig:topomap-5band}). The signal is carried by an
\emph{anger-versus-warm-positive} contrast on $9$ of the $28$ video
clips: cohort $r{=}+0.87$ at PO3/$\gamma$ for Anger, Amusement, and
Tenderness, against $r{=}-0.02$ on the remaining $19$ mid-valence
stimuli. Removing Anger alone shrinks the cohort correlation to
$+0.33$. Full per-region tables, the Simpson's-paradox per-subject
caveat, time-resolved LPP-window dynamics, and theta--gamma
phase-amplitude coupling tests are in
Appendix~\ref{app:topography-deep-dive}.

\begin{figure}[ht]
\centering
\IfFileExists{figures/neuro/NF1_5band_topomaps.pdf}
  {\includegraphics[width=\linewidth]{neuro/NF1_5band_topomaps.pdf}}
  {\includegraphics[width=\linewidth]{F4_brain_topography.pdf}}
\caption[Brain topography]{\textbf{V-axis encoding peaks in posterior visual cortex across all
five frequency bands.}
\textbf{(a--e)} MNE-style scalp topomaps, one per band ($\delta$,
$\theta$, $\alpha$, $\beta$, $\gamma$), of the cohort Pearson
correlation between the LLM V-axis ($28$ stimulus projections) and
per-channel band power on FACED ($n{=}123$ subjects). Diverging
colourmap: positive (red) vs.\ negative (blue) correlation; symmetric
limits across bands. Gold circles mark the top channel per band;
text under each topomap names that channel and its $r$ value
(strongest cell PO3/$\gamma$, $r{=}{+}0.48$). The colourbar (centred)
gives Pearson $r$; the small head diagram (right) shows the
$32$-channel 10--20 montage used in all five panels (A1/A2 mastoids
excluded from plotting). \textbf{(f)} Region-mean $|r|$ aggregated
over all five bands: occipital ($0.21$) $>$ parietal ($0.18$) $>$
central ($0.18$) $>$ frontal ($0.16$). The cohort signal is
posterior-dominant for video-evoked emotion---a paradigm-level
finding for dynamic stimuli rather than a refutation of the static
frontal-alpha asymmetry literature, which we replicate qualitatively
in direction at smaller magnitude.
}
\label{fig:topomap-5band}
\end{figure}

\paragraph{This is not a refutation of frontal-alpha asymmetry.}
Davidson's frontal-alpha asymmetry
hypothesis~\citep{davidson1992anterior,coan2004frontal} was developed
on static stimuli; the cohort signal we report is on dynamic
$28$-second video. Right-minus-left frontal-alpha asymmetry remains
positive in our data ($+0.006$ to $+0.016$ across F-pairs),
consistent with FAA's predicted direction at smaller magnitude than
the posterior $|r|$. This is a paradigm-level finding: in the
dynamic-video regime the cohort signal lives where the visual content
does, not where Davidson's theory of static-image
approach/withdrawal predicts. We do not claim the V-axis subsumes FAA
or that posterior visual cortex is the seat of valence; we claim
that, on dynamic video, the cohort signal that aligns with the
LLM-derived direction is carried by visual-content processing of
high-arousal clips. Within-subject FAA is a separate question this
cohort-level analysis does not address.
